\documentclass[14pt]{beamer}
\usepackage{pepslides}
\title{Equilibrium figures}
\lecturenumber{22}
\hypersetup{pdftitle={Lecture 22: Equilibrium figures}}

\begin{document}

\titleslide

\begin{frame}[t]
\heading{From last time\ldots}
Lagrangian in a self-gravitating planet:
\begin{eqn}
\mathcal{L}(\mathbf{x},t,\bphi,\mathbf{v},\mathbf{F})=\frac{1}{2}\rho(\mathbf{x})\|\mathbf{v}\|^{2}
-W(\mathbf{x},t,\mathbf{F})-\frac{1}{2}\rho(\mathbf{x})\zeta(\mathbf{x},t).
\end{eqn}
Decomposition of the motion relative to an approximately rotating frame:
\begin{eqn}
\varphi_{i}(\mathbf{x},t)=\Phi_{i}(t)+R_{ij}(t)\overline{\varphi}_{j}(\mathbf{x},t).
\end{eqn}
Equilibrium equations:
\begin{eqn}
\rho\epsilon_{ijk}\epsilon_{klm}\Omega_{j}\Omega_{l}x_{m}-\frac{\partial T_{ij}^{0}}{\partial x_{j}}-\rho\gamma_{i}^{0}=0.
\end{eqn}
\end{frame}

\begin{frame}
\centering
\includegraphics[height=0.72\paperheight]{chandrasekhar1969_newton_cassini.png}
\credit{Chandrasekhar (1969)}
\end{frame}

\sectionslide{Solutions of the equilibrium equations}

\begin{frame}
Notation from last lecture:
\begin{eqn}
\rho\epsilon_{ijk}\epsilon_{klm}\Omega_{j}\Omega_{l}x_{m}-\frac{\partial T_{ij}^{0}}{\partial x_{j}}-\rho\gamma_{i}^{0}=0.
\end{eqn}
To simplify things we now write:
\begin{eqn}
\rho\epsilon_{ijk}\epsilon_{klm}\Omega_{j}\Omega_{l}x_{m}-\frac{\partial\sigma_{ij}}{\partial x_{j}}-\rho g_{i}=0.
\end{eqn}
At equilibrium, referential and spatial quantities coincide.

The stress tensor is symmetric and so has six independent components.

The force balance holds in general materials, not just elastic ones.
\end{frame}

\begin{frame}
Suppose the planet's shape, density and angular velocity are known.

We then have \textbf{three} equations for the \textbf{six} components of stress:
\begin{eqn}
\rho\epsilon_{ijk}\epsilon_{klm}\Omega_{j}\Omega_{l}x_{m}-\frac{\partial\sigma_{ij}}{\partial x_{j}}-\rho g_{i}=0.
\end{eqn}
This is an underdetermined system and so we \textbf{cannot expect unique solutions}.

More formally, there exist \textbf{non-trivial divergence-free symmetric second-order tensor fields} that satisfy
the correct boundary conditions.

Compare with the existence of divergence-free vector fields.
\end{frame}

\sectionslide{Hydrostatic equilibrium}

\begin{frame}
The (Cauchy) stress in a compressible Newtonian fluid:
\begin{eqn}
\sigma_{ij}=-p\,\delta_{ij}+\eta\left(\frac{\partial v_{i}}{\partial x_{j}}+\frac{\partial v_{j}}{\partial x_{i}}-\frac{2}{3}\frac{\partial v_{k}}{\partial x_{k}}\delta_{ij}\right).
\end{eqn}
At equilibrium, the stress is \textbf{hydrostatic}:
\begin{eqn}
\sigma_{ij}=-p\,\delta_{ij}.
\end{eqn}
The same result holds in more general viscous fluids.

Scaling estimate for deviatoric stress associated with convection:
\begin{eqn}
\eta V/L\sim10^{6}\,\mathrm{Pa}.
\end{eqn}
Pressures in the Earth's interior are of order $10^{10}\,\mathrm{Pa}$ -- deviation from hydrostatic state in
Earth-like planets represents only a \textbf{small perturbation}.
\end{frame}

\begin{frame}
Given the hydrostatic approximation, we have:
\begin{eqn}
\rho\epsilon_{ijk}\epsilon_{klm}\Omega_{j}\Omega_{l}x_{m}+\frac{\partial p}{\partial x_{i}}-\rho g_{i}=0.
\end{eqn}
This is now an \textbf{overdetermined} set of equations -- three equations for one unknown.

Helpful to switch to a vector-calculus notation:
\begin{eqn}
\nabla p+\rho\nabla\gamma=\mathbf{0},\qquad\gamma=\phi+\psi.
\end{eqn}
Here we have introduced the \textbf{gravity potential}, this being the \textbf{sum of the gravitational
and centrifugal potentials}.
\end{frame}

\begin{frame}
Hydrostatic equilibrium equation:
\begin{eqn}
\nabla p+\rho\nabla\gamma=\mathbf{0}.
\end{eqn}
Gradients of pressure and gravity potential are parallel, and so a level surface of one field is a level surface
of the other.

Ignoring the atmosphere, the pressure vanishes at the planet's surface, and hence this must be a level surface
of the gravity potential.

Taking the curl of the equation, we find:
\begin{eqn}
\nabla\rho\times\nabla\gamma=\mathbf{0}.
\end{eqn}
This shows that level surfaces of density are also coincident with the other two fields.
\end{frame}

\begin{frame}
\centering
\includegraphics[height=0.6\paperheight]{L10F1.pdf}

\raggedright
In hydrostatic equilibrium the level surfaces of $p$, $\rho$ and $\gamma$ all coincide, and the surface is
an equipotential of gravity.
\end{frame}

\begin{frame}
\begin{columns}[c]
\begin{column}{0.5\textwidth}
\centering
\includegraphics[width=\textwidth]{myhill2025_phobos_potential.png}
\end{column}
\begin{column}{0.45\textwidth}
Equatorial slice through the gravitational potential for Phobos calculated assuming a constant interior density.

The potential is not constant on the moon's surface, and hence this cannot be a hydrostatic equilibrium state.

In general, the hydrostatic equilibrium condition imposes very strong conditions on a planet's shape, density,
and rotation.
\end{column}
\end{columns}
\credit{Myhill \emph{et al.} (2025)}
\end{frame}

\begin{frame}
In the non-rotating case, suppose the planet is spherically symmetric.

Equilibrium equation reduces to:
\begin{eqn}
\nabla p+\rho\frac{\ddns\phi}{\ddns r}\hat{\mathbf{r}}=\mathbf{0}.
\end{eqn}
There is just one non-trivial equation here, and so we can find an equilibrium solution:
\begin{eqn}
\frac{\ddns p}{\ddns r}+\rho\frac{\ddns\phi}{\ddns r}=0,\qquad p(b)=0.
\end{eqn}
It can be shown that this is the \textbf{only} solution in the non-rotating case.
\end{frame}

\begin{frame}
By definition, a planet is \textbf{slowly rotating} if its centrifugal potential is small relative to its
gravitational potential:
\begin{eqn}
\frac{\textrm{centrifugal potential}}{\textrm{gravitational potential}}\sim\frac{3\Omega^{2}}{8\pi G\overline{\rho}}.
\end{eqn}
For the Earth, we find:
\begin{eqn}
\frac{\textrm{centrifugal potential}}{\textrm{gravitational potential}}\sim2\times10^{-3}.
\end{eqn}
In such cases, we can apply first-order perturbation theory.

Higher-order theories exist and are sometimes necessary.
\end{frame}

\begin{frame}
Expand about the non-rotating solution, with $s$ a perturbation parameter:
\begin{eqn}
\begin{aligned}
\rho&=\rho_{0}+s\rho_{1}+\cdots,\\
p&=p_{0}+s\,p_{1}+\cdots,\\
\phi&=\phi_{0}+s\,\phi_{1}+\cdots,\\
\psi&=s\,\psi_{1}.
\end{aligned}
\end{eqn}
Similarly, the free-surface radius is expanded as:
\begin{eqn}
r(\theta,\varphi)=b+s\,h_{1}(\theta,\varphi)+\cdots.
\end{eqn}
Topography on internal boundaries is expanded similarly, but we neglect this detail for simplicity in what follows.
\end{frame}

\begin{frame}
First-order perturbation of the equilibrium equation:
\begin{eqn}
\nabla p_{1}+\rho_{0}\nabla\gamma_{1}+\rho_{1}g_{0}\hat{\mathbf{r}}=\mathbf{0},
\end{eqn}
with $\gamma_{1}=\phi_{1}+\psi_{1}$ and $g_{0}=\ddns\phi_{0}/\ddns r$.

Exact boundary condition on the perturbed planetary surface:
\begin{eqn}
p[b+s\,h_{1}(\theta,\varphi)+\cdots,\theta,\varphi]=0.
\end{eqn}
Expanded to first order, and using $\ddns p_{0}/\ddns r=-\rho_{0}g_{0}$, this reduces to:
\begin{eqn}
\begin{gathered}
p_{1}+\frac{\ddns p_{0}}{\ddns r}h_{1}=0,\\[\eqgap]
p_{1}-\rho_{0}g_{0}h_{1}=0.
\end{gathered}
\end{eqn}
\end{frame}

\begin{frame}
First-order perturbation of the equilibrium equation:
\begin{eqn}
\nabla p_{1}+\rho_{0}\nabla\gamma_{1}+\rho_{1}g_{0}\hat{\mathbf{r}}=\mathbf{0}.
\end{eqn}
Take the cross product with the radial unit vector:
\begin{eqn}
\begin{gathered}
\hat{\mathbf{r}}\times\nabla p_{1}+\hat{\mathbf{r}}\times\rho_{0}\nabla\gamma_{1}=\mathbf{0},\\[\eqgap]
\hat{\mathbf{r}}\times\nabla(p_{1}+\rho_{0}\gamma_{1})=\mathbf{0}.
\end{gathered}
\end{eqn}
Taking the perturbations to average to zero over spherical surfaces, we conclude that:
\begin{eqn}
p_{1}=-\rho_{0}\gamma_{1}.
\end{eqn}
\end{frame}

\begin{frame}
First-order perturbation of the equilibrium equation:
\begin{eqn}
\nabla p_{1}+\rho_{0}\nabla\gamma_{1}+\rho_{1}g_{0}\hat{\mathbf{r}}=\mathbf{0}.
\end{eqn}
Taking the curl of this equation, we arrive at:
\begin{eqn}
\hat{\mathbf{r}}\times\nabla\left(\frac{\ddns\rho_{0}}{\ddns r}\gamma_{1}-g_{0}\rho_{1}\right)=\mathbf{0}.
\end{eqn}
And hence, we see that:
\begin{eqn}
\rho_{1}=\frac{1}{g_{0}}\frac{\ddns\rho_{0}}{\ddns r}\gamma_{1}.
\end{eqn}
\end{frame}

\begin{frame}
We have now shown that the following identities hold for the first-order perturbations:
\begin{eqn}
p_{1}=-\rho_{0}\gamma_{1},\quad h_{1}=-\frac{\gamma_{1}}{g_{0}},\quad\rho_{1}=\frac{1}{g_{0}}\frac{\ddns\rho_{0}}{\ddns r}\gamma_{1}.
\end{eqn}
All other terms can be expressed through the gravity potential perturbation.
\end{frame}

\begin{frame}
Linearisation of Poisson's equation is trivial:
\begin{eqn}
\nabla^{2}\phi_{1}=\begin{cases} 4\pi G\rho_{1} & r<b, \\ 0 & r>b. \end{cases}
\end{eqn}
Linearisation of the continuity conditions requires more work:
\begin{eqn}
[\phi_{1}]_{-}^{+}=0,\quad[\hat{\mathbf{r}}\cdot\nabla\phi_{1}+4\pi G\rho_{0}h_{1}]_{-}^{+}=0.
\end{eqn}
Derivation given in an appendix to the notes (non-examinable).
\end{frame}

\begin{frame}
Expressing other perturbations in terms of the gravity potential:
\begin{eqn}
\nabla^{2}\phi_{1}=\begin{cases} 4\pi G\dfrac{1}{g_{0}}\dfrac{\ddns\rho_{0}}{\ddns r}(\phi_{1}+\psi_{1}) & r<b, \\ 0 & r>b, \end{cases}
\end{eqn}
along with the continuity conditions:
\begin{eqn}
[\phi_{1}]_{-}^{+}=0,\quad\left[\hat{\mathbf{r}}\cdot\nabla\phi_{1}-4\pi G\rho_{0}\frac{(\phi_{1}+\psi_{1})}{g_{0}}\right]_{-}^{+}=0.
\end{eqn}
Linear boundary value problem for the perturbed gravitational potential, with the centrifugal potential providing
the ``force terms''.

Theory applies also to equilibrium tides in gaseous planets.
\end{frame}

\begin{frame}
The centrifugal potential is proportional to $Y_{20}$, and separation of variables shows that:
\begin{eqn}
\phi_{1}+\psi_{1}=\sqrt{\frac{4\pi}{5}}\frac{2}{3}rg_{0}\epsilon\,Y_{20},
\end{eqn}
which serves to define the \textbf{ellipticity}, $\epsilon(r)$. Figures of slowly rotating planets are oblate spheroids.

At the planet's mean surface (and at internal ones too) we have:
\begin{eqn}
\epsilon(b)=\frac{\textrm{equatorial radius}-\textrm{polar radius}}{\textrm{mean radius}}.
\end{eqn}
Surface ellipticity of the Earth is around 1/300.
\end{frame}

\begin{frame}
The ellipticity satisfies \textbf{Clairaut's equation} (Clairaut 1743):
\begin{eqn}
\frac{\ddns^{2}\epsilon}{\ddns r^{2}}+8\pi G\rho_{0}g_{0}^{-1}\left(\frac{\ddns\epsilon}{\ddns r}+\frac{\epsilon}{r}\right)-\frac{6\epsilon}{r^{2}}=0,
\end{eqn}
subject to the boundary conditions:
\begin{eqn}
\left.\frac{\ddns\epsilon}{\ddns r}\right|_{r=0}=0,\quad\left.\frac{\ddns\epsilon}{\ddns r}\right|_{r=b}=\frac{1}{b}\left[\frac{5\Omega^{2}b^{3}}{2GM}-2\epsilon(b)\right].
\end{eqn}
Solved numerically by integrating outwards from $r=0$ and rescaling to meet the condition at $r=b$.
\end{frame}

\begin{frame}
\centering
\includegraphics[height=0.55\paperheight]{L10F2.pdf}

\raggedright
Left: radial density profile obtained from analysis of the Earth's free oscillations -- final two lectures.

Right: the corresponding ellipticity (or flattening), plotted here as a reciprocal.

Hydrostatic surface value 1/299.9, observed 1/298.3: the difference reflects non-hydrostatic stresses in the mantle.
\end{frame}

\begin{frame}[t]
\heading{What is a planet?}
\begin{columns}[b]
\begin{column}{0.45\textwidth}
\centering
\includegraphics[height=0.5\paperheight]{nasa_newhorizons_pluto.png}

A planet
\end{column}
\begin{column}{0.45\textwidth}
\centering
\includegraphics[width=0.9\textwidth]{nasa_hirise_phobos.png}

Not a planet
\end{column}
\end{columns}
\credit{NASA}
\end{frame}

\sectionslide{Rapidly rotating bodies (non-examinable)}

\begin{frame}
\centering
\includegraphics[height=0.66\paperheight]{L10F3.pdf}

\raggedright
Homogeneous fluid at any rotation rate: \textbf{Maclaurin spheroids} (1742) and \textbf{Jacobi ellipsoids} (1834).

The Jacobi sequence \textbf{bifurcates} from the Maclaurin sequence; pear-shaped figures branch in turn from the Jacobi sequence.
\end{frame}

\begin{frame}
\centering
\begin{tabular}{cc}
\raisebox{-0.5\height}{\includegraphics[height=0.42\paperheight]{haumea_render.jpg}} &
\raisebox{-0.5\height}{\includegraphics[height=0.48\paperheight]{L10F4.pdf}}
\end{tabular}

\raggedright
Equilibrium figures \textbf{do not have to be axisymmetric}.

Haumea is a dwarf planet with a rotation period of 3.9 hours and a triaxial figure, its axes within 1\,\% of a
Jacobi ellipsoid (Ortiz \emph{et al.} 2017).

It rotates too fast to be a homogeneous fluid, and so must have a denser core.
\end{frame}

\begin{frame}
\centering
\includegraphics[height=0.7\paperheight]{darwin_fission_illustration.png}

\raggedright
\textbf{Fission hypothesis} for the Moon's formation, due to George Darwin. Long dismissed in favour of a giant
impact, but it does have a real physical basis, and it has recently attracted renewed interest.
\end{frame}

\begin{frame}[t]
\heading{What you need to know and be able to do}
\begin{enumerate}
\item[(i)] Know the form of the equilibrium equations, and that they hold in a general planet and not just an elastic
one.
\item[(ii)] That the equilibrium equations are underdetermined, and why.
\item[(iii)] What a hydrostatic equilibrium state is, and be able to derive the constraints imposed on the level
surfaces of $p$, $\rho$, and $\gamma$.
\item[(iv)] Derive the hydrostatic equilibrium state in a non-rotating planet.
\item[(v)] Know in outline how solutions can be obtained within a slowly rotating planet. Here you will not need to
reproduce the full derivation, but may be asked to perform related calculations with appropriate guidance.
\end{enumerate}
\end{frame}

\end{document}
